% !TEX program = xelatex
% =====================================================================
% 全国大学生数学建模竞赛（CUMCM）论文主文件
% 来源：Mrite 项目 format.cls（基于 MIT 许可的 cumcmthesis 改造，内嵌思源宋体）
% 编译：xelatex -interaction=nonstopmode 论文.tex （连续两遍解决交叉引用）
% 说明：本模板使用 withoutpreface 选项跳过承诺书页；如需承诺书页去掉该选项。
% =====================================================================
\PassOptionsToPackage{quiet}{xeCJK}
\documentclass[withoutpreface,bwprint]{format}
\usepackage{etoolbox}
\usepackage{ctex}
\BeforeBeginEnvironment{tabular}{\zihao{-5}}
\usepackage[framemethod=TikZ]{mdframed}
\usepackage{url}
\usepackage{array}
\usepackage{tabularx}
\usepackage{longtable}
\newcolumntype{C}{>{\centering\arraybackslash}X}
\newcolumntype{R}{>{\raggedleft\arraybackslash}X}
\newcolumntype{L}{>{\raggedright\arraybackslash}X}

\renewcommand{\textbf}[1]{{\song\bfseries #1}}

\usepackage{tikz}
\usetikzlibrary{arrows.meta}
\tikzset{
  box/.style={rectangle, draw, minimum width=3.2cm, minimum height=0.9cm, align=center},
  arrow/.style={thick, -{Stealth}}
}

\title{论文标题}

\begin{document}

{\centering\zihao{3}\bfseries 论文标题\par}

\setcounter{page}{1}
\thispagestyle{empty}

% ===================== 摘要 =====================
% 规则（详见技能 references/paper-writing.md）：
%   总字数 <= 900 字且必须一页；开头段 <= 30 字；结尾段 <= 20 字。
%   每问一段，按问题类型弹性分配字数：建模求解型 30-35%、数据分析型 20-25%、建议总结型 10-15%。
%   每问必须出现具体数值结果；关键词 5 个左右。
\begin{abstract}

% 开头段（<=30 字）：
...是...领域的重要研究课题。本文基于...数据，围绕...展开系统研究。

% 逐问摘要（每问选一种类型模板替换【】）：
% ① 有数据型：
\textbf{针对问题一，}本问是【类型】问题，需【要求】。首先对附件数据预处理，提取【特征】，【清洗/标准化】。接着进行【变量筛选/降维/参数选择】。然后建立【模型】进行【求解过程】。最后通过【对比/验证】，结果表明\textbf{【关键数值结果】}，模型有效地【评价】。

% ② 无数据（机理）型：
\textbf{针对问题二，}本问是【类型】问题，需【要求】。首先分析【机理/物理过程/目标约束】，建立【方程/优化模型】。接着设定【参数/初始条件】。然后通过【数值方法/算法】求解，得到【结果】。结果表明\textbf{【关键数值结果】}，模型有效地【评价】。

% ③ 建议总结型（一笔带过）：
\textbf{针对问题三，}本问是建议总结问题。基于前文分析，从【角度1】和【角度2】出发给出【N】条建议，核心结论为\textbf{【要点】}。

% 结尾段（<=20 字）：
本文模型性能优良、结论可靠，可为【领域】提供决策参考。

\keywords{关键词1；关键词2；关键词3；关键词4；关键词5}
\label{abstract:end}
\end{abstract}

\thispagestyle{empty}

\tableofcontents
\thispagestyle{empty}
\newpage

\setcounter{page}{1}

% ===================== 一、引言 =====================
\section{引言}

\subsection{问题背景}
% 三段式：宏观背景 -> 问题来源与困境 -> 解决渴求。自然段落，不分点。

\subsection{问题重述}
% 每问 <= 2 行，删背景只留问题本质；不使用分点符号。
\textbf{问题1：}...

\textbf{问题2：}...

\textbf{问题3：}...

% ===================== 二、总体分析 =====================
% 三段式纯文字，不画流程图：
%   第1段：一句话总体表达全文思路；
%   第2段：针对问题X建立了XX模型来解决XX（逐问串联）；
%   第3段：总结各问题之间的递进关系。
\section{总体分析}

...

% ===================== 三、模型假设 =====================
% 两段式：引导句 + itemize。假设条数 = 问题数（或按实际需要），
% 每条 <= 2 行，且每条假设都应在正文建模处被实际引用。
\section{模型假设}

为简化问题，本文做出以下假设：

\begin{itemize}[itemindent=2em]
\item \textbf{假设1：}...
\item \textbf{假设2：}...
\item \textbf{假设3：}...
\end{itemize}

% ===================== 四、符号说明 =====================
% 全论文所有表格统一使用 longtable 样式（见下），禁用 tabular/tabularx。
% 列宽通式（按标准版心 455pt 校准，预留 \tabcolsep 每列 12pt）：
%   各列 p{比例\textwidth} 的比例总和 = 1 - 0.03 * 列数N
%   2列=0.94, 3列=0.91, 4列=0.88, 5列=0.85, 6列=0.82, 7列=0.79, 8列=0.76 ...
% 所有列必须带 >{\centering\arraybackslash}，否则报 Misplaced \noalign。
\section{符号说明}

本文使用的主要符号及其含义如表\ref{tab:符号说明}所示。

\begin{longtable}{>{\centering\arraybackslash}p{0.16\textwidth}>{\centering\arraybackslash}p{0.78\textwidth}}
  \caption{符号说明}
  \label{tab:符号说明} \\
  \toprule
  符号 & 说明 \\
  \midrule
  \endfirsthead
  \caption{符号说明（续）} \\
  \toprule
  符号 & 说明 \\
  \midrule
  \endhead
  \bottomrule
  \endfoot
  $x$ & ... \\
  $y$ & ... \\
  $z_{ij}$ & ... \\
\end{longtable}

% ===================== 五、模型的建立与求解（主控文件） =====================
% 按题目实际问题数 N 动态调整：N>4 新增 \input，N<4 删除多余行。
\section{模型的建立与求解}

% ===================== 5.1 问题1的建立求解（装配文件） =====================
% 每问复制本文件改名（5.2、5.3...），内部固定拆为两个子文件。
% 拆分原则：模型建立与模型求解必须在同一装配下（求解大量引用建立的公式与变量）。
\subsection{问题一模型的建立和求解}

% ===================== 5.1.1 分析与准备 =====================
% 结构：具体分析（三段文字 + 固定9框流程图） + 模型准备（数据预处理 + 算法选择）
% 禁止：讲检验、讲结果（此时尚未求解）。

\subsubsection{问题一的具体分析}
% 第1段：针对问题X + 问题类型 + 大致方法；
% 第2段：首先...接着...然后...最后...串联完整步骤，不换段；
% 第3段：具体分析流程如图\ref{fig:analysis1_flow}所示。

\begin{figure}[ht]
  \begin{center}
    \begin{tikzpicture}[x=1cm, y=0.7cm]
      \node[box] (n11) at (0, 0) {...};
      \node[box] (n12) at (5, 0) {...};
      \node[box] (n13) at (10, 0) {...};
      \node[box] (n22) at (5, -3) {...};
      \node[box] (n21) at (0, -3) {...};
      \node[box] (n23) at (10, -3) {...};
      \node[box] (n31) at (0, -6) {...};
      \node[box] (n32) at (5, -6) {...};
      \node[box] (n33) at (10, -6) {...};
      \draw[arrow] (n11) -- (n12);
      \draw[arrow] (n12) -- (n13);
      \draw[arrow] (n23) -- (n22);
      \draw[arrow] (n22) -- (n21);
      \draw[arrow] (n13) -- ++(0,-1.0) -| (n23);
      \draw[arrow] (n21) -- (n31);
      \draw[arrow] (n31) -- (n32);
      \draw[arrow] (n32) -- (n33);
    \end{tikzpicture}
    \caption{问题一分析流程图}
    \label{fig:analysis1_flow}
  \end{center}
\end{figure}

\subsubsection{数据预处理}
% 仅问题1写本小节（数据预处理五段结构）：
%   第1段 预处理原因（缺失/量纲/异常）；第2段 数据理解与提取（特征数、样本规模）；
%   第3段 预处理步骤（清洗->标准化->特征衍生，公式写出）；
%   第4段 预处理结果（处理前后 min/max/mean/std 对比）；第5段 总结。
% Min-Max 归一化与 Z-score 标准化公式：
% x' = (x - x_min)/(x_max - x_min)      x' = (x - mu)/sigma

\subsubsection{算法选择}
% 问题1必写。五段式：定性 -> 候选算法与评估维度 -> longtable 对比表（优/良/中/差，
% 必须写具体算法名，禁止"算法A"）-> 对比分析 -> 总结引出。
% 评估维度按问题类型调整（预测类可用：预测精度/训练效率/可解释性/鲁棒性）。

\begin{longtable}{>{\centering\arraybackslash}p{0.17\textwidth}>{\centering\arraybackslash}p{0.17\textwidth}>{\centering\arraybackslash}p{0.17\textwidth}>{\centering\arraybackslash}p{0.17\textwidth}>{\centering\arraybackslash}p{0.17\textwidth}}
  \caption{算法能力对比}
  \label{tab:算法对比1} \\
  \toprule
  算法名称 & 维度一 & 维度二 & 维度三 & 维度四 \\
  \midrule
  \endfirsthead
  \caption{算法能力对比（续）} \\
  \toprule
  算法名称 & 维度一 & 维度二 & 维度三 & 维度四 \\
  \midrule
  \endhead
  \bottomrule
  \endfoot
  遗传算法（GA） & 优 & 优 & 良 & 良 \\
  蚁群算法（ACO） & 良 & 中 & 优 & 中 \\
  模拟退火（SA） & 中 & 良 & 良 & 中 \\
\end{longtable}


% ===================== 5.1.2 建模与求解 =====================
% 模型建立（三段式，占最大篇幅）：
%   第1段（引文）：针对问题X，本节将从XX角度建立XX模型...具体过程如下。
%   第2段（展开）：按 \textbf{步骤N：标题} 编号展开，每步 6-7 个公式，
%     公式间用文字串联（"将式(X)代入式(Y)"），每式后紧跟变量说明段
%     （"其中，$X$表示...，$Y$表示..."，自然段不分点）。
%   第3段（桥接）：综上所述，本节完成了XX模型的建立，接下来进行求解。
% 四类问题的编号步骤骨架（详见 references/paper-writing.md）：
%   优化类：  1.目标函数的建立 -> 2.约束条件的建立 -> 3.核心机理建模 -> 4.求解算法建模
%   评价类：  1.评价指标体系构建 -> 2.指标权重确定（熵权/AHP 完整推导） -> 3.综合评价方法
%   预测/分类：1.特征工程 -> 2.模型选定与构建 -> 3.模型训练与评估
%   机理类：  1.机理分析框架 -> 2.关键变量关系建模 -> 3.模型整合与参数确定
% 模型求解：引文 -> 表/图 -> 分析。结果数值用 longtable，趋势对比用图（图宽 0.8\textwidth）。

\subsubsection{模型的建立}

\textbf{步骤1：}...描述。
\begin{equation}
  f(x) = ...
\end{equation}
其中，$x$表示...，$f(x)$表示...。

\textbf{步骤2：}...描述。
\begin{equation}
  g(x) = ...
\end{equation}
其中，...。

\subsubsection{模型的求解}

...求解过程与结果分析，关键数值以 longtable 呈现，并引用每张图表。


% \input{5.2.问题2的建立求解.tex}
% \input{5.3.问题3的建立求解.tex}

% ===================== 六、模型检验与分析 =====================
% 结构：引文 -> 二级标题分类 -> 文字 -> longtable -> 分析。
% 子节：误差分析（优先 5 折交叉验证）+ 灵敏度分析（关键参数 ±20% 扰动）。
% 灵敏度数据在求解阶段前置产出，此处直接引用。
\section{模型检验与分析}

为验证本文构建的模型的有效性和稳健性，进行以下检验。

\subsection{误差分析}

\begin{longtable}{>{\centering\arraybackslash}p{0.22\textwidth}>{\centering\arraybackslash}p{0.22\textwidth}>{\centering\arraybackslash}p{0.22\textwidth}>{\centering\arraybackslash}p{0.22\textwidth}}
  \caption{模型误差分析结果}
  \label{tab:error} \\
  \toprule
  指标 & 训练集 & 测试集 & 说明 \\
  \midrule
  \endfirsthead
  \caption{模型误差分析结果（续）} \\
  \toprule
  指标 & 训练集 & 测试集 & 说明 \\
  \midrule
  \endhead
  \bottomrule
  \endfoot
  $R^2$ & ... & ... & 决定系数 \\
  RMSE & ... & ... & 均方根误差 \\
  MAE & ... & ... & 平均绝对误差 \\
  MAPE(\%) & ... & ... & 平均绝对百分比误差 \\
\end{longtable}

如表\ref{tab:error}所示，...

\subsection{灵敏度分析}

\begin{longtable}{>{\centering\arraybackslash}p{0.30\textwidth}>{\centering\arraybackslash}p{0.30\textwidth}>{\centering\arraybackslash}p{0.30\textwidth}}
  \caption{灵敏度分析结果}
  \label{tab:sensitivity} \\
  \toprule
  参数变化 & 输出变化 & 灵敏度评价 \\
  \midrule
  \endfirsthead
  \caption{灵敏度分析结果（续）} \\
  \toprule
  参数变化 & 输出变化 & 灵敏度评价 \\
  \midrule
  \endhead
  \bottomrule
  \endfoot
  ... & ... & ... \\
\end{longtable}

如表\ref{tab:sensitivity}所示，...

% ===================== 七、模型评价 =====================
% 统一规则（融合裁决）：优点 4 条 + 缺点 2 条（2-3 条亦可），优点必须多于缺点。
% 按模型写，不按问题写；每条 <= 3 行；每条都要有具体依据，不写空话。
\section{模型的评价}

\subsection{模型的优点}
\begin{itemize}[itemindent=2em]
\item \textbf{优点1：}...
\item \textbf{优点2：}...
\item \textbf{优点3：}...
\item \textbf{优点4：}...
\end{itemize}

\subsection{模型的缺点}
\begin{itemize}[itemindent=2em]
\item \textbf{缺点1：}...
\item \textbf{缺点2：}...
\end{itemize}

% ===================== 八、模型的改进与推广 =====================
% 各一段自然段落，不分点。
\section{模型的改进与推广}

\subsection{模型的改进}
...

\subsection{模型的推广}
...


\newpage
% ===================== 九、参考文献 =====================
% 规则：
%   1. 必须根据论文实际内容生成，禁止保留示例文献；只写真实存在的文献。
%   2. 数量 8-15 条；GB/T 7714-2015 格式；纯文本一行写完；
%      禁用 \emph/\textbf/\textit/\underline 等格式命令；删除 DOI。
%   3. 标签：作者姓氏+年份+关键词，如 \bibitem{behzadian2012topsis}。
%   4. 每条必须在正文被 \cite{} 引用；写完后逐一核对一一对应。
%   5. 引用位置：模型建立（方法原始文献）、模型求解（算法文献）、
%      模型检验（对比方法文献）、引言（领域背景文献）。
\begin{thebibliography}{99}

\bibitem{ref1} ...

\bibitem{ref2} ...

\end{thebibliography}


\newpage
% ===================== 十、附录 =====================
% 附件说明表：统一 longtable 两列。不写运行环境/硬件环境。
\makeatletter
\let\@oldaddcontentsline\addcontentsline
\renewcommand{\addcontentsline}[3]{}
\section*{附录}
\let\addcontentsline\@oldaddcontentsline
\makeatother

\begin{longtable}{>{\raggedright\arraybackslash}p{0.38\textwidth}>{\centering\arraybackslash}p{0.56\textwidth}}
  \caption{附件说明表}
  \label{tab:appendix} \\
  \toprule
  文件 & 说明 \\
  \midrule
  \endfirsthead
  \caption{附件说明表（续）} \\
  \toprule
  文件 & 说明 \\
  \midrule
  \endhead
  \bottomrule
  \endfoot
  问题一\_求解.py & 问题1求解代码 \\
  预处理数据.csv & 预处理后的数据 \\
  ... & ... \\
\end{longtable}


\end{document}
